% !TeX root = ../thuthesis-example.tex

\chapter{相关技术概述}

\section{Spring Boot后端开发技术}

Spring Boot是Java生态中常用的后端开发框架，它通过自动配置、起步依赖和内嵌Web容器降低了服务搭建成本。对于本项目而言，Spring Boot主要承担REST接口、业务编排、数据访问、鉴权拦截和推荐算法调度等职责。系统可按照Controller、Service、Repository、Domain、Config等层次组织代码，使接口入口、业务规则和数据库访问边界清晰。

在推荐系统场景中，Spring Boot的优势体现在三个方面。第一，接口层可以使用\texttt{@RestController}和\texttt{@RequestMapping}快速定义推荐Feed、趋势内容、用户推荐和行为埋点接口。第二，服务层可以通过依赖注入组合多个召回源、过滤器和评分器，便于替换算法策略。第三，配合Spring Security和JWT可以对用户身份进行校验，避免推荐接口在未授权情况下暴露用户画像和行为数据。

\section{Spring Cloud与分布式思想}

本项目虽然以课程设计为主，但后端设计遵循Spring Cloud的分布式思想。推荐系统可以拆分为用户服务、群组服务、内容服务、推荐服务、实验服务和日志服务。实际部署时，这些服务既可以先以单体Spring Boot应用中的多模块方式实现，也可以在业务增长后逐步拆分为微服务。

在分布式系统中，推荐服务需要关注调用链稳定性和缓存策略。用户画像、关注关系和内容特征来自不同数据源，若每次请求都同步查询全部数据，会导致接口响应变慢。因此系统将部分热点数据放入Redis缓存，并对召回源设置超时时间。当某一路召回失败时，系统可以回退到热门内容或冷启动内容，保证Feed接口可用。

\section{Supabase PostgreSQL云端数据库}

Supabase是基于PostgreSQL的后端即服务平台，提供托管数据库、认证、存储和实时能力。本项目使用Supabase作为云端数据库，项目名称为\texttt{telegram}，区域为\texttt{us-west-2}，数据库版本为PostgreSQL 17。通过MCP读取可知，数据库中已经存在用户、群组、联系人、新闻内容、用户行为、用户兴趣向量和实验配置等表。

相比本地数据库，Supabase的优势是部署简单、云端访问方便、适合课程设计展示。对于推荐系统而言，\texttt{news\_articles}表保存内容特征，\texttt{news\_user\_events}表保存点击、停留和分享等行为，\texttt{news\_user\_vectors}表保存短期和长期兴趣向量，\texttt{experiments}表保存实验分桶策略。这些表共同支撑个性化推荐的输入数据。

\section{Redis缓存技术}

Redis是内存型键值数据库，常用于缓存、排行榜、计数器和分布式锁。本项目在推荐系统中使用Redis承担三个职责：

\begin{enumerate}
  \item 时间线缓存。将作者最新内容写入有序集合，例如按作者维护最近七天内容，用户请求时按关注列表聚合候选内容。
  \item 热点内容缓存。将高点击、高分享和高互动内容缓存为热门候选，作为冷启动和召回失败时的兜底来源。
  \item 临时特征缓存。缓存用户近期行为、已曝光内容和请求追踪标识，减少数据库重复查询。
\end{enumerate}

Redis缓存不替代数据库持久化，而是作为提高吞吐量和降低延迟的中间层。推荐接口每次请求都可能触发多路召回和特征补全，缓存可以显著减少数据库压力。

\section{Go消息投递与Redis Stream消费技术}

Go语言具有轻量级协程、通道和高效网络库，适合实现高并发消息消费、异步投递和后台任务执行。本项目在Spring Boot业务服务之外，引入Go侧消息投递服务，主要处理平台事件、消息投递、同步唤醒和outbox状态更新。该服务不直接替代业务后端，而是承担高频、异步、可恢复的投递链路，使主业务接口可以将耗时的多接收者分发工作转移到后台执行。

Go投递服务的核心输入来自Redis Stream。Redis Stream支持消费者组、消息确认、待处理消息列表（PEL）和消息重放，适合表达“至少一次投递”和失败恢复语义。为了避免全局串行消费带来的吞吐瓶颈，Go侧按照\texttt{chat\_id}优先、其次\texttt{user\_id}、\texttt{outbox\_id}或\texttt{message\_id}的方式计算消息分区键，将消息派发到多个worker lane。同一业务键在同一lane内保持顺序，不同业务键可以并行处理，从而兼顾即时通信系统中的顺序要求和大群消息扇出的并发需求。

在确认机制上，服务将成功消息按stream聚合后批量\texttt{XACK}，减少Redis远程往返；失败消息仍保留原有重试、恢复和死信语义。对于PEL reclaim，系统将其从主消费循环中拆分为独立的有预算后台worker，避免旧消息恢复阻塞新消息读取。为了降低MongoDB写放大，投递链路还引入分块update id预留、outbox聚合状态增量维护和批量wake事件发布。上述设计体现了分布式消息系统中常见的三个原则：顺序只在必要的业务键内保证、远程操作尽量批处理、恢复逻辑不阻塞主路径。

\section{C++图计算与快照查询技术}

C++适合实现对延迟和内存布局敏感的图查询内核。本项目在推荐召回阶段需要频繁查询用户的邻居关系、共同联系人、共同群组和多跳社交路径。如果所有图关系都在通用业务服务中逐条查询数据库，会造成较高的远程访问成本和不可控的尾延迟。因此系统引入C++图计算服务，将社交图谱预构建为内存快照，通过HTTP接口向推荐服务提供高性能图查询能力。

图计算服务采用快照发布模型。后端服务负责提供图快照数据，C++服务定期从\texttt{/internal/graph-kernel/snapshot}拉取并构建内存索引。快照构建阶段将用户id和边类型intern为整数编号，再按源点组织为紧凑CSR（Compressed Sparse Row）结构。CSR结构使用连续数组保存邻居和offset，相比散列表或对象数组可以减少内存碎片，提高缓存命中率，也便于对高出度节点执行顺序扫描和TopK截断。

在查询层，C++服务支持直接邻居查询、ranked邻居查询、overlap共同关系查询和多跳遍历。为了控制高出度节点带来的开销，系统将请求级时间\texttt{now\_ms}采样一次，并在ranking和traversal中复用同一个时间基准，避免热循环重复系统调用。overlap查询采用streaming top-K思想，不必先物化全部交集再排序，而是在扫描过程中维护候选上限。多跳遍历采用budget-aware策略，在访问预算内优先保留更有价值的候选。服务还为快照版本、快照表示、构建阶段耗时、扫描数、访问数和预算耗尽情况输出summary和metrics，方便推荐服务判断图召回是否健康。

\section{Apifox接口联调技术}

Apifox是一体化API协作工具，支持接口设计、请求调试、接口文档和自动化测试。本项目使用Apifox对Spring Boot后端接口进行联调。测试时可以在请求栏中输入本地服务地址，例如\texttt{http://localhost:8080/api/space/feed}，并通过Params区域配置\texttt{limit}、\texttt{in\_network\_only}等参数。

在课程报告中，Apifox截图不仅用于展示测试工具，也用于证明接口路径、请求方法、参数拆解和联调过程。后续如果接入真实Spring Boot运行服务，还可以在Apifox中保存接口集合，设置环境变量、Authorization请求头和响应断言。

\section{前端React与TypeScript技术}

系统前端采用React与TypeScript思想进行设计。React负责组件化页面展示，TypeScript为接口数据结构提供类型约束。前端主要页面包括用户登录、群组与联系人、新闻列表、推荐Feed、趋势内容、用户推荐和行为反馈。前端向后端发送HTTP请求，后端返回JSON格式数据，前端再根据返回的内容渲染列表、卡片和交互按钮。

在推荐系统中，前端不仅展示推荐结果，也承担行为采集入口。例如用户点击新闻、停留阅读、分享内容或屏蔽某条内容时，前端可以向后端行为埋点接口上报事件。后端将事件写入\texttt{news\_user\_events}，推荐服务再根据这些行为更新兴趣向量和排序权重。

\section{个性化推荐算法技术}

个性化推荐算法的核心任务是从大量候选内容中选出最适合当前用户的一组内容。本项目采用多阶段推荐架构，主要包括召回、过滤、排序和重排四个阶段。

召回阶段负责扩大候选集。系统可以从关注作者内容、社交图谱扩展内容、新闻向量近邻、作者兴趣匹配、热门内容和冷启动内容中分别取出候选。不同召回源适合不同用户状态：新用户更依赖热门内容和冷启动策略，活跃用户更适合使用兴趣向量和历史行为，关注关系较丰富的用户则适合使用社交图谱。

过滤阶段负责保证候选内容可用。系统会去除重复内容、用户自己发布的内容、已经曝光过的内容、来自屏蔽用户的内容以及命中敏感词的内容。过滤阶段强调规则确定性，优先保证安全和体验。

排序阶段负责计算候选内容得分。本项目采用加权融合思想，核心特征包括点击、回复、转发、分享、停留时长、作者亲和度、内容质量、新鲜度和负反馈。正向行为提升得分，负向行为降低得分。例如回复和转发代表较强互动意愿，权重高于普通点击；屏蔽、举报和“不感兴趣”属于强负反馈，需要显著降低相关内容排序。

重排阶段负责控制最终结果的多样性。TopK选择不能只选最高分内容，否则可能出现同一作者、同一主题或同一来源连续占据列表的情况。系统在选择阶段加入作者多样性、主题多样性和内容形态约束，使推荐结果更加均衡。
